{
\renewcommand{\arraystretch}{1.5}
\setlist[enumerate]{leftmargin=*, labelsep=0.5em}

\begin{tabularx}{\textwidth}{|>{\raggedright\arraybackslash}X|>{\raggedright\arraybackslash}X|}
\hline
\textbf{Tên Use case} & Xóa bài viết \\
\hline
\textbf{ID} & UC-0039 \\
\hline
\textbf{Tác nhân} & Người dùng đã xác thực (tác giả bài viết) \\
\hline
\textbf{Mô tả} & Người dùng xóa bài viết blog đã đăng \\
\hline
\textbf{Điều kiện tiên quyết} & Bài viết tồn tại, người dùng là tác giả \\
\hline
\textbf{Điều kiện sau} & Bài viết bị xóa khỏi hệ thống \\
\hline
\textbf{Kích hoạt} & Người dùng chọn xóa bài viết \\
\hline
\textbf{Luồng chính} &
\begin{enumerate}
    \item Người dùng mở bài viết của mình.
    \item Người dùng nhấn Xóa.
    \item Hệ thống hiển thị cảnh báo xác nhận.
    \item Người dùng xác nhận xóa.
    \item Hệ thống xóa bài viết khỏi cơ sở dữ liệu.
    \item Người dùng được chuyển hướng về danh sách bài viết.
\end{enumerate} \\
\hline
\textbf{Luồng thay thế} &
\begin{enumerate}
    \item \textbf{Quản trị viên xóa bài viết:}
    \begin{enumerate}
        \item[1A.] Quản trị viên có thể xóa bất kỳ bài viết nào bất kể quyền sở hữu.
    \end{enumerate}
\end{enumerate} \\
\hline
\textbf{Luồng ngoại lệ} &
\begin{enumerate}
    \item \textbf{Bài viết không tồn tại:}
    \begin{enumerate}
        \item[1Err.] Hệ thống trả về lỗi 404.
    \end{enumerate}
    \item \textbf{Không có quyền:}
    \begin{enumerate}
        \item[1Err2.] Hệ thống trả về lỗi 403 khi người dùng không phải tác giả hoặc Admin.
    \end{enumerate}
\end{enumerate} \\
\hline
\end{tabularx}
\captionof{table}{Đặc tả Use case: UC-0039 Xóa bài viết}
}
